% ----------------------------------------------------------------------------------
% Chapter 4: System Architecture and Implementation
% ----------------------------------------------------------------------------------
% Detail the hardware (LIMO Cobot, MyCobot, Jetson Orin Nano) and the software (ROS 2 nodes, topics, actions).
\label{ch:implementation}

\section{Hardware Platform}
The Agilex LIMO Cobot integrates a mobile base with...

\section{Software Architecture in ROS 2}
The software stack is built sequentially using independent, communicating nodes...

\section{Simulation Environment (Gazebo)}
The system is developed within a simulation-to-reality (sim-to-real) workflow in which Gazebo
Classic constitutes the first phase rather than the end goal. The complete navigation,
perception, and manipulation pipeline is built and validated in simulation---where errors carry
no hardware cost, the scenario is fully observable, and the world can be reset
deterministically---before the same ROS~2 system is deployed onto the physical LIMO and
myCobot~280 platform available in the laboratory. The simulated robot is constructed to mirror
the real platform's kinematics, sensor placement, and coordinate frames as closely as is
practical, so that the transition to hardware alters only the lowest layer of the stack: the
Gazebo physics and sensor plugins are replaced by the corresponding hardware drivers, while the
perception, motion-planning, and orchestration nodes are reused unchanged. This separation of a
hardware-agnostic upper stack from a substitutable hardware layer is the principal engineering
rationale for prioritising a high-fidelity simulation.

\section{Perception Subsystem}
Object localisation is performed by a deliberately lightweight, classical
computer-vision pipeline rather than a learned detector, on the grounds of
transparency, the absence of training data, and ease of diagnosis for a single
known-coloured target. The depth camera publishes an RGB image, a registered depth
image, and the camera intrinsics. The target is segmented in the HSV colour space
(robust to illumination relative to raw RGB), and the centroid pixel $(u,v)$ is
extracted. The corresponding depth $Z$ is read from the depth image and the pixel is
de-projected into a metric point in the camera frame using the pin-hole model,
$X = (u-c_x)Z/f_x$, $Y = (v-c_y)Z/f_y$. This point is transformed into the
manipulator planning frame (\texttt{base\_link}) via the TF tree and published as a
\texttt{PoseStamped} on \texttt{/box\_pose}. A documented limitation of the layout is
that the manipulator's reachable workspace lies partly inside the camera's near-clip
distance, so the reachable grasp region and the camera-visible region only narrowly
overlap; this is analysed quantitatively in Chapter~\ref{ch:experiments}.

\section{Motion Planning with MoveIt~2}
Arm motion is planned with MoveIt~2. As the Python bindings (\texttt{moveit\_py}) are
not packaged for the target distribution (ROS~2 Humble), the orchestration node
communicates with the running \texttt{move\_group} through the standard
\texttt{moveit\_msgs} action and service interfaces over \texttt{rclpy}, requiring no
additional dependencies. The semantic model (SRDF) defines the six-joint planning
group, named configurations, and the link-pair collision matrix; inverse kinematics
is solved by the KDL plugin and paths are searched by OMPL's RRTConnect. Each grasp
sub-goal (pre-grasp, grasp, lift, place) is issued as a Cartesian pose goal; the
planner computes a collision-checked joint trajectory from the current state, which is
executed by the \texttt{ros2\_control} joint-trajectory controller. The gripper is
driven outside MoveIt by direct position commands, with ramped set-points to avoid the
impulsive contact forces that a step command induces.

\section{Engineering Challenge: An Incomplete Collision Matrix}
A significant integration fault is documented here because its diagnosis illustrates a
class of failure common to MoveIt deployments. Joint-space goals executed correctly,
yet \emph{every} Cartesian goal was rejected and the arm folded to its home
configuration. Direct interrogation of the inverse-kinematics service confirmed the
targets were reachable, and the state-validity service confirmed both start and goal
configurations were collision-free; planning nonetheless failed whenever the waist
joint deviated from zero. Examination of the planner log and interpolation of the
candidate path through the state-validity checker revealed a phantom self-collision:
the auto-generated collision matrix disabled checks between the gripper \emph{fingers}
and the mobile base, but omitted the wrist body links (\texttt{gripper\_base},
\texttt{joint6}, \texttt{joint6\_flange}). Consequently, any forward-facing
configuration was incorrectly reported as colliding with the base wheels, steering,
laser, and camera, and was refused. Restoring the omitted \texttt{disable\_collisions}
entries resolved all Cartesian planning. The episode underlines that an incomplete
semantic collision model silently invalidates motion planning while leaving joint-space
control unaffected, and that disciplined, service-level verification is required to
localise the cause.

\section{Simulation Physical Fidelity and Stability}
Obtaining numerically stable behaviour in Gazebo Classic required several specific
measures, each addressing a distinct failure mode. First, every link's inertia tensor
must satisfy the rigid-body triangle inequalities, $I_{xx}+I_{yy}\ge I_{zz}$ and its
cyclic permutations; a manufacturer-supplied tensor that violates this condition causes
the ODE solver to evaluate an undefined state and the model to diverge immediately on
load. Second, rigid-body contact between a position-controlled gripper and a support
surface generates large separating impulses that destabilise a comparatively light and
unconstrained base; this was mitigated by terminating descents clear of supporting
surfaces, registering those surfaces as planning-scene collision objects, ramping
gripper set-points to avoid impulsive commands, constraining the base pose during
manipulation, and---most directly---softening the surface contact model by lowering its
stiffness and capping the ODE contact-correction velocity (\texttt{max\_vel}), so that
any residual penetration is resolved gently rather than as an explosive impulse. Third,
the mobile base exhibits uncommanded creep under solver noise and
manipulator reaction forces unless the drive joints specify explicit friction and
damping. Fourth, because Gazebo Classic does not enforce URDF mimic constraints, the
passive joints of the parallel gripper must be driven through the \texttt{ros2\_control}
hardware interface to remain kinematically consistent. A related fault was a
half-configured gripper in which the controller commanded a single master joint while
the finger joints were declared as both mimic and command interfaces; this contradictory
specification left the gripper inert. It was resolved by driving all six gripper joints
directly with mirror-signed set-points, which is robust to Gazebo Classic's unreliable
mimic handling. These measures are characteristic of manipulator simulation in Gazebo
Classic and are reported as part of the methodology.

\section{Model Fidelity and Parameter Verification}
Because the simulation is the first phase of a sim-to-real workflow, the physical
parameters of the model must correspond to the real platform; an unverified value
silently introduces a discrepancy that surfaces only on hardware. A systematic audit was
therefore performed, tracing every physically meaningful parameter in the URDF and
sensor descriptions to a manufacturer specification, or else explicitly designating it
as a simulation-tuning or test-scenario value. The platform was confirmed from the
procurement record to be the AgileX LIMO~PRO with an Elephant Robotics myCobot~280, an
Orbbec DaBai depth camera, an EAI~T-mini~Pro LiDAR, and an NVIDIA Jetson Orin~Nano.
Cross-checking the model against these specifications revealed and corrected several
inherited errors: the wheel track ($0.14\rightarrow0.172$\,m, confirmed from the
mechanical datasheet), the depth-camera
horizontal field of view ($80^{\circ}\rightarrow67.9^{\circ}$) and depth range and
resolution ($0\text{--}8\rightarrow0.3\text{--}3$\,m, $640\times480\rightarrow640\times400$),
and the LiDAR field of view and range
($240^{\circ},\,0.2\text{--}8\,\text{m}\rightarrow360^{\circ},\,0.02\text{--}12\,\text{m}$).
The adaptive gripper's clamping range was confirmed from its datasheet as $20\text{--}45$\,mm,
informing the choice of a $30$\,mm manipulation target.

The audit also surfaced the most consequential sim-to-real discrepancy. The physical
Orbbec DaBai has a minimum depth of $0.30$\,m, whereas the simulated sensor's near-clip
had been lowered to $0.05$\,m so that a close target remained visible. An object at the
$\approx0.21$\,m grasp distance is therefore visible in simulation but lies
\emph{inside the blind spot of the physical camera}. This is retained deliberately and
flagged: the physical system must perceive the target from a greater stand-off distance
and approach under a remembered pose, rather than perceive-and-grasp in place. The
governing principle adopted is that no physical parameter may remain an unjustified
constant: each must trace to a specification or be tagged as tuning or test.

\section{Tool-Frame Calibration and Grasp Geometry}
Reliable grasping required calibrating the manipulator's tool frame, a step whose
necessity is itself a finding. Three geometric corrections were identified. First, the
planning frame \texttt{gripper\_tcp} is not coincident with the centre of the finger
gap, so aiming the planning frame at the target left the fingers offset by approximately
one centimetre; this was corrected with a constant lateral grasp offset. Second, the
pose-goal orientation tolerance, initially $0.5$\,rad ($\approx29^{\circ}$), was large
enough to admit visibly tilted ``top-down'' solutions; tightening it to $0.1$\,rad
produced a vertical gripper. Third, the adaptive gripper's fingers translate downward as
they close, so the closed tool reaches lower than the open tool; the grasp height was
raised so that the closing fingers clear the support surface. These corrections---lateral
offset, orientation tolerance, and closing clearance---constitute the tool-frame
calibration that, on the physical robot, generalises to a combined hand--eye and
tool-centre-point calibration.

\section{Navigation Integration and Drive Architecture}
The differential-drive kinematics of the LIMO PRO is exposed to Nav2 via the
\texttt{gz\_diff\_drive} Gazebo plugin and the \texttt{diff\_drive\_controller} from
\texttt{ros2\_control}. The measured wheel track of the LIMO PRO is 172\,mm; this
value was confirmed from the manufacturer's mechanical datasheet and set as the
\texttt{wheel\_separation} parameter in the drive controller. A discrepancy of 3\,mm
from an earlier assumed value of 175\,mm is small in absolute terms but accumulates
over distance during dead reckoning, making accurate track width an important parameter
for odometry consistency with the physical robot.

The simulation odometry source is set to \texttt{WORLD}, which instructs Gazebo to
derive the odometric pose directly from the simulated world frame rather than by
integrating wheel velocities. The consequence is that the \texttt{odom} frame in the
TF tree is numerically identical to the Gazebo world frame: pose errors and accumulated
drift are absent. This arrangement is intentional: in the simulation validation phase
the objective is to verify that the planning, perception, and manipulation layers are
individually correct before introducing the compounding uncertainty of real odometry.
The simulation navigation pipeline therefore constitutes an upper bound on achievable
performance, against which the hardware trials---using real odometry---will be compared
as part of the sim-to-real evaluation.

A Nav2 configuration file (\texttt{nav2\_limo\_diff.yaml}) parameterises the local and
global costmaps, the AMCL particle filter, the DWB local planner, and the behaviour
trees governing goal approach and recovery. The most significant departure from a stock
configuration is the costmap obstacle minimum range. In a stock Nav2 deployment the
minimum range at which LiDAR returns are admitted into the costmap is set to 0.22\,m.
For the LIMO Cobot this value causes the arm, which extends 0.25--0.35\,m in the
rest posture, to be treated as an obstacle, marking the robot's own vicinity as a
lethal zone and rendering the local planner unable to find any valid trajectory. This
was diagnosed by observing the local costmap in RViz: a ring of lethal cells
surrounded the robot even in an open environment. The fix is to raise both
\texttt{obstacle\_min\_range} and \texttt{raytrace\_min\_range} in both costmaps to
0.35\,m, which places the near-field arm geometry below the sensor admission threshold.

\section{Sensor Configuration: LiDAR Self-Hit and SLAM Map Quality}
An important implementation fault, and its correction, is documented here because it
demonstrates a category of error that is specific to robots where the sensor's nominal
field of view subtends the robot's own body.

The EAI T-mini Pro LiDAR fitted to the LIMO PRO has a nominal $360^{\circ}$ horizontal
field of view. When the sensor model was initially parameterised with this specification,
beams in the angular range beyond $\pm108^{\circ}$ from the forward axis struck the
robot's own chassis sides at approximately 6--7\,cm range. The sensor minimum range was
set to 0.02\,m, which is less than this chassis return distance, so these hits were
accepted as valid obstacle measurements. During SLAM mapping the robot treated its own
body as a cluster of moving obstacles, baking 55 spurious occupied cells into the map.
When Nav2 used this contaminated map for global planning and AMCL for localisation,
the phantom walls degraded path quality and destabilised scan-match convergence.

The correction applied two parameters simultaneously. The LiDAR field of view was
restricted to $\pm120^{\circ}$ (a 2.094\,rad half-angle), which is conservative
with respect to the chassis shadow angle ($\pm108^{\circ}$) and eliminates all chassis
returns by geometry. The minimum valid range was raised from 0.02\,m to 0.10\,m as
a defence in depth, rejecting any residual near-field return regardless of angle.
A new SLAM session with the corrected sensor produced a map of 238$\times$358 cells
(at 0.05\,m/cell) containing no spurious obstacle pixels in open traversable space,
with 97.2\,\% of cells classified as free.

The AMCL minimum range parameter was set to match the sensor threshold (0.10\,m) to
ensure that the particle filter's expected measurement model agreed with the observed
data; a mismatch between the sensor minimum range and the AMCL model range causes
degraded likelihood scoring on near-range returns and reduces localisation accuracy.

\section{Self-Collision Policy: Manipulator versus Mobile Base}
A deliberate modelling decision concerns collision checking between the manipulator and
the mobile-base body. The two subsystems are modelled at different collision fidelities:
the arm links carry accurate convex meshes, whereas the base and sensors are represented
by coarse primitive volumes. The base collision primitive extends into the volume the
folded arm occupies, so enabling arm--body collision checking produces false-positive
(phantom) collisions even in postures where the physical parts do not touch; empirically
this rendered \emph{every} configuration, including the home posture, unplannable.
Arm--body checks are therefore disabled in simulation, while manipulator self-collision
(arm against arm) remains fully enabled, so the dominant hazard---the arm folding into
its own links---is still prevented. This is safe with respect to the physical robot,
which carries accurate factory collision geometry and the manufacturer's semantic model,
in which arm--body collision is checked correctly. The disable is thus a
simulation-modelling accommodation for coarse meshes, not a removal of real-robot
safety; the corresponding refinement---tighter base meshes and regeneration of the
collision matrix---is identified as future work.

\section{Integrated Navigation--Manipulation Orchestration}
The complete task cycle---navigate to the pickup location, pick the target object,
back up---is coordinated by a single ROS~2 node, \texttt{nav\_pick\_orchestrator}.
The node holds action clients for Nav2's \texttt{NavigateToPose} and MoveIt's
\texttt{MoveGroup} interfaces, and a service client for the custom
\texttt{grasp\_attacher} weld service; it sequences these interfaces in a deterministic
state machine without a dedicated behaviour-tree engine. The sequencing is as follows.

\begin{enumerate}
    \item \textbf{Transport posture.} Before commanding any base motion, the node waits
    for the MoveIt \texttt{move\_group} action server to accept connections and then
    commands the arm to the \texttt{travel} configuration
    ($\theta = [0, 1.2, {-0.6}, {-0.6}, 0, 0]$\,rad). This folds the arm downward,
    lowering the overall centre of mass and shortening the lever arm of the second joint,
    which is a horizontal-axis hinge whose compliance causes it to act as an inverted
    pendulum under base acceleration. Without a compact transport posture, inertial
    forces during navigation excite oscillations at joint~2 that persist for several
    seconds after the base stops.
    \item \textbf{Navigation.} A \texttt{NavigateToPose} goal is sent to Nav2 with the
    pre-configured table waypoint. The node blocks until Nav2 reports success or failure.
    \item \textbf{Visual docking.} On arrival at the coarse navigation goal, the node
    engages a visual-servo docking stage: an HSV-thresholded centroid from the RGB
    camera drives proportional linear velocity commands until the estimated box range
    falls within the arm's grasp envelope.
    \item \textbf{Grasp and retract.} The arm is moved to the \texttt{ready}
    configuration, the box 3D pose is requested from the perception node, and the
    standard grasp sequence (pre-grasp $\to$ descent $\to$ close $\to$ lift) is
    executed. On closure the weld service rigidly attaches the box to the gripper link.
    \item \textbf{Backup.} The base reverses by a fixed displacement to clear the
    pickup zone.
\end{enumerate}

The node architecture makes a deliberate simplification: the \texttt{odom} frame, which
in the Gazebo simulation is identical to the world frame (\texttt{odometry\_source:
WORLD}), is used as the common reference for all TF lookups during manipulation.
Specifically, the grasp-attacher service resolves the wrist position as a live TF lookup
(\texttt{odom}$\to$\texttt{gripper\_tcp}) at the moment of weld, rather than propagating
a stored pre-navigation position. This is essential: when the base moves from its startup
position to the pickup table, any TF offset stored before navigation is invalidated; a
live lookup at the instant of grasping is always consistent regardless of the path
taken by the base.

\section{Simulation Rendering: Dual-GPU Segfault and NVIDIA PRIME Offload}
A hardware-specific fault that is documented here because it affects reproducibility on
consumer laptops commonly available in university laboratories. The development machine
is equipped with both a discrete NVIDIA RTX~3060 Mobile GPU and an AMD Cezanne
integrated GPU, configured in Linux PRIME on-demand mode. In this mode, applications
render on the integrated GPU by default; the discrete GPU receives workloads only when
explicitly requested via environment variables.

Gazebo Classic's GPU sensor plugin (used by the depth camera) executed on the AMD
integrated GPU in this configuration. The Gazebo depth-camera sensor consistently
caused a segmentation fault (exit code $-11$, SIGSEGV) in \texttt{gzserver} during or
immediately after model spawn, before any control loop had started. The fault occurred
reliably under the default PRIME policy and did not occur under software rendering
(\texttt{LIBGL\_ALWAYS\_SOFTWARE=1}), confirming the root cause as the OpenGL driver
stack used by the integrated GPU rather than anything in the ROS or Gazebo configuration.

Software rendering was used as a temporary workaround: it eliminated the crash at the
cost of CPU-side rendering of the depth camera ($640\times400$ at 10\,Hz), which reduced
the simulation's real-time factor to well below 1.0 and made interactive development
slow. The permanent fix was to wire NVIDIA PRIME offload environment variables
(\texttt{DRI\_PRIME=1}, \texttt{\_\_NV\_PRIME\_RENDER\_OFFLOAD=1},
\texttt{\_\_GLX\_VENDOR\_LIBRARY\_NAME=nvidia}) into the Gazebo launch file's
\texttt{additional\_env} dictionary, so that both \texttt{gzserver} and \texttt{gzclient}
render on the RTX~3060 by default. This restored full GPU-accelerated rendering, with
both processes confirmed via \texttt{nvidia-smi} to consume GPU memory on the discrete
card (147\,MiB and 142\,MiB respectively), and the depth camera continued to publish
correctly without any driver crash. No change to the ROS or Gazebo configuration was
required beyond the environment variables in the launch file.

This class of PRIME-related Gazebo crash is a known issue on Optimus/PRIME laptops under
Ubuntu with proprietary NVIDIA drivers and is not specific to this platform; the fix
described here is portable to any PRIME-configured machine.

\section{RViz Configuration and MoveIt Panel Incompatibility}
The default RViz configuration used for navigation monitoring (\texttt{nav\_pick.rviz})
included a MoveIt \texttt{MotionPlanning} display panel. On the same dual-GPU machine,
loading this configuration caused RViz to segfault during startup. The crash was
consistent and hardware-specific: the MoveIt panel's robot model loading triggered the
same AMD OpenGL driver issue identified in the previous section, but in this case the
PRIME offload variables did not prevent it because RViz was launched as a child process
that did not inherit the Gazebo launch environment.

The fix was to create a dedicated navigation monitoring configuration
(\texttt{nav\_monitor.rviz}) that is identical to \texttt{nav\_pick.rviz} except that
the \texttt{moveit\_rviz\_plugin/MotionPlanning} display block is absent. The navigation
monitoring view contains: RobotModel, TF tree, LaserScan, Map, global and local
costmaps, the Nav2 planned path, and the RGB camera feed. MoveIt trajectory
visualisation, when needed, is performed in a separate RViz session launched from within
the MoveIt pipeline, where the PRIME environment is correctly propagated. This separation
of navigation and manipulation visualisation contexts resolved the crash and is retained
as a deliberate architectural choice.

\section{Gazebo World Physics and Platform Stability}
The Gazebo ODE physics engine's contact-correction velocity parameter
(\texttt{contact\_max\_correcting\_vel}) governs the speed at which interpenetrating
bodies are pushed apart after each collision step. A world file inherited from an
upstream reference repository set this parameter to 0.1\,m/s, compared to the
Gazebo default of 100\,m/s. At 0.1\,m/s the contact solver is effectively unable
to rapidly resolve penetrations in high-CoM configurations: when the arm extends
forward, the manipulator mass raises the platform's effective centre of mass well
above the wheel contact points; even a small solver disturbance then accumulates over
multiple time steps rather than being corrected, eventually destabilising the base
and causing it to tip forward. The parameter was removed, reverting to the Gazebo
default (100\,m/s), which resolves contacts within a single integration step and
yields stable simulation over the full navigation and manipulation cycle.
